\chapter{Теоретические основы применения ML и LLM в малом и среднем бизнесе}

\section{Роль данных в управлении МСП}

Для малого и среднего бизнеса данные имеют прикладную ценность только тогда, когда они помогают принимать решения быстрее и точнее. В производственной компании такими решениями являются закупка сырья, планирование выпуска, контроль возвратов, управление платежным календарем, оценка зависимости от клиентов и поставщиков.

В отличие от крупного бизнеса, МСП часто не имеет полноценного хранилища данных и выделенной команды аналитиков. Поэтому практический ML-контур должен быть простым, воспроизводимым и интерпретируемым. Для такой компании важны не только точность модели, но и понятность признаков, возможность ручной проверки и устойчивость к изменению форматов выгрузок.

\section{Типовые задачи машинного обучения}

В рамках финансового и операционного анализа малого производственного бизнеса можно выделить несколько групп задач:
\begin{itemize}
  \item классификация банковских операций по управленческим категориям;
  \item прогноз спроса и заказов;
  \item прогноз возвратов и корректировок;
  \item прогноз дневного cash flow;
  \item выявление аномальных поступлений и платежей;
  \item мониторинг концентрации клиентов и поставщиков.
\end{itemize}

Классификация операций нужна для того, чтобы банковская выписка стала управленческим отчетом. Прогнозные модели помогают заранее увидеть риск кассового разрыва или перепроизводства. Аномалии полезны как очередь операций на проверку, а не как автоматическое решение о блокировке платежа.

\section{Rule-based, ML и LLM-подходы}

Rule-based подход основан на заранее заданных правилах: ключевых словах, контрагентах, типах операций и регулярных паттернах. Его преимущество -- прозрачность и контроль. Недостаток -- трудоемкая поддержка правил при появлении новых контрагентов или формулировок назначений платежа.

Классические ML-модели применимы для прогнозирования спроса, возвратов и денежных потоков. При короткой истории данных особенно полезны модели с лаговыми и календарными признаками, например gradient boosting или CatBoost. Они способны использовать общие закономерности между товарами, контрагентами и периодами.

LLM-подход полезен для интерпретации текстовых назначений платежей, особенно когда формулировки неоднородны. Однако в финансовой задаче LLM не должна быть единственным источником истины: результат нужно ограничивать строгим списком категорий, проверять JSON-ответ и сравнивать с rule-based разметкой.

\section{Ограничения применения ML в МСП}

Основные ограничения связаны с качеством и объемом данных. История наблюдений может составлять всего 12--14 месяцев, контрагенты могут иметь разные названия в банке и 1С, часть справочников может быть неполной, а крупные разовые платежи способны искажать прогноз.

Поэтому для МСП предпочтителен гибридный подход: правила и ручная валидация используются для контроля качества, ML -- для прогноза и обнаружения закономерностей, LLM -- для поддержки текстовой классификации.
